% !TEX program = pdflatex
\documentclass[aspectratio=169,10pt]{beamer}

% =========================
% PAQUETES
% =========================
\usepackage[spanish]{babel}
\usepackage[T1]{fontenc}
\usepackage[utf8]{inputenc}
\usepackage{lmodern}
\usepackage{amsmath,amssymb,array}
\usepackage{booktabs}
\usepackage{xcolor}

% =========================
% COLORES AGRONOMÍA (UST)
% =========================
\definecolor{USTGreen}{HTML}{2E7D32}
\definecolor{USTLightGreen}{HTML}{A5D6A7}
\definecolor{USTDark}{HTML}{1B5E20}
\definecolor{USTGray}{HTML}{2F2F2F}

\setbeamercolor{normal text}{fg=USTGray,bg=white}
\setbeamercolor{structure}{fg=USTGreen}
\setbeamercolor{frametitle}{fg=white,bg=USTGreen}
\setbeamercolor{block title}{bg=USTGreen, fg=white}
\setbeamercolor{block body}{bg=USTLightGreen!20, fg=black}

\setbeamercolor{exampleblock title}{bg=USTDark, fg=white}
\setbeamercolor{exampleblock body}{bg=USTLightGreen!10, fg=black}

\setbeamercolor{alertblock title}{bg=red!70!black, fg=white}
\setbeamercolor{alertblock body}{bg=red!10, fg=black}

\setbeamertemplate{navigation symbols}{}

% =========================
% DATOS
% =========================
\title{Fundamentos de Matemática Básica}
\subtitle{Unidades de Medida en Agronomía}
\author{Profesor Luis Tulio González}
\institute{Universidad Santo Tomás \\ Departamento de Ciencias Básicas}
\date{Semestre I - 2026}

% =========================
\begin{document}
	
	% =========================
	\begin{frame}
		\titlepage
	\end{frame}
	
	% =========================
	\begin{frame}{Objetivos de la clase}
		\begin{itemize}[<+->]
			\item Comprender el concepto de magnitud.
			\item Identificar unidades de medida en contextos agrícolas.
			\item Convertir unidades correctamente.
			\item Resolver problemas aplicados a la agronomía.
			\item Detectar errores e inconsistencias en cálculos.
		\end{itemize}
	\end{frame}
	
	% =========================
	\begin{frame}{¿Qué es medir?}
		\begin{block}{Idea clave}
			Medir es comparar una cantidad con una unidad de referencia.
		\end{block}
		
		\begin{itemize}[<+->]
			\item Superficie de terreno
			\item Cantidad de fertilizante
			\item Volumen de agua
			\item Tiempo de riego
		\end{itemize}
		
		\pause
		\begin{exampleblock}{Ejemplo}
			\(1500\ \text{L de agua}\)
		\end{exampleblock}
	\end{frame}
	
	% =========================
	\begin{frame}{Magnitudes}
		\begin{block}{Definición}
			Una magnitud es toda propiedad que puede medirse.
		\end{block}
		
		\begin{itemize}[<+->]
			\item Longitud
			\item Masa
			\item Tiempo
			\item Área
			\item Volumen
		\end{itemize}
		
		\pause
		\begin{alertblock}{Importancia}
			Permiten tomar decisiones correctas en producción agrícola.
		\end{alertblock}
	\end{frame}
	
	% =========================
	\begin{frame}{Longitud}
		\[
		km \quad hm \quad dam \quad m \quad dm \quad cm \quad mm
		\]
		
		\begin{itemize}[<+->]
			\item → derecha: $\times 10$
			\item ← izquierda: $\div 10$
		\end{itemize}
		
		\pause
		\begin{exampleblock}{Ejemplo}
			\(0.4\ m = 40\ cm\)
		\end{exampleblock}
	\end{frame}
	
	% =========================
	\begin{frame}{Masa}
		\[
		kg \quad hg \quad dag \quad g \quad dg \quad cg \quad mg
		\]
		
		\begin{itemize}[<+->]
			\item → derecha: $\times 10$
			\item ← izquierda: $\div 10$
		\end{itemize}
		
		\pause
		\begin{exampleblock}{Ejemplo}
			\(250\ g = 0.25\ kg\)
		\end{exampleblock}
	\end{frame}
	
	% =========================
	\begin{frame}{Área}
		\begin{block}{Clave}
			Cada salto = $\times 100$
		\end{block}
		
		\[
		m^2
		\]
		
		\pause
		\[
		1\ ha = 10\,000\ m^2
		\]
		
		\pause
		\begin{exampleblock}
			\(2\ ha = 20\,000\ m^2\)
		\end{exampleblock}
	\end{frame}
	
	% =========================
	\begin{frame}{Volumen}
		\begin{block}{Clave}
			Cada salto = $\times 1000$
		\end{block}
		
		\pause
		\[
		1\ L = 1000\ mL
		\]
		
		\pause
		\begin{exampleblock}
			\(2.5\ L = 2500\ mL\)
		\end{exampleblock}
	\end{frame}
	
	% =========================
	\begin{frame}{Ejemplo aplicado}
		\begin{block}{Fertilizante}
			300 g por planta, 80 plantas
		\end{block}
		
		\pause
		\[
		300 \times 80 = 24\,000\ g
		\]
		
		\pause
		\[
		24\,000\ g = 24\ kg
		\]
		
		\pause
		\begin{exampleblock}
			Se requieren 24 kg
		\end{exampleblock}
	\end{frame}
	
	% =========================
	\begin{frame}{Errores comunes}
		\begin{alertblock}{Cuidado}
			\begin{itemize}[<+->]
				\item Sumar distintas unidades
				\item Confundir ha con m$^2$
				\item No convertir antes de operar
			\end{itemize}
		\end{alertblock}
	\end{frame}
	
	% =========================
	\begin{frame}{Ejercicio}
		\begin{itemize}[<+->]
			\item \(4.2\ kg \rightarrow ?\ g\)
			\item \(1.8\ ha \rightarrow ?\ m^2\)
			\item \(3.5\ L \rightarrow ?\ mL\)
		\end{itemize}
	\end{frame}
	
	% =========================
	\begin{frame}{Cierre}
		\begin{itemize}[<+->]
			\item Medir = comparar
			\item Las unidades importan
			\item Convertir evita errores
			\item Aplicación directa en agronomía
		\end{itemize}
	\end{frame}
	
	% =========================
	\begin{frame}
		\centering
		\Huge ¿Preguntas?
	\end{frame}
	
\end{document}